% ==========================================
% CHAPTER 08: Level 10 (Google Authenticator)
% ==========================================
\chapter{Level 10: The Cryptographic TOTP Masterclass}

\textbf{In This Chapter:}
\begin{itemize}
    \item Systemically analyzing Single-Point-of-Failure architectures.
    \item Utilizing standard Python `hmac` implementations in MicroPython.
    \item Utilizing the `api.qrserver.com` endpoint for QR generation over Wi-Fi.
    \item Full code execution for NodeMCU Demo 10 and Pico W Authenticator.
\end{itemize}

Our device is impenetrable to automated scripts, but we harbor a catastrophic \textbf{Single-Point-Of-Failure}. If an intern accidentally posts the password `1234` into a public GitHub repository, or a technician types it on a sticky note attached to the server rack, the lockout defense becomes completely bypassed because the algorithm trusts the keyholder implicitly. 

\section{The TOTP Mechanics (\url{https://datatracker.ietf.org/doc/html/rfc6238})}
Zero-Trust infrastructure mandates validating intersecting factors. The modern implementation standard for "Possession" validation is the \textbf{Google Authenticator} app on a smartphone.

1. Both the IoT device and the phone share an identical cryptographic base variable (e.g., `JBSWY3DP`) permanently stored in isolating memory constraints.
2. Both independent hardwares synchronize with a universal atomic clock via UDP Network Time Protocol vectors.
3. Both silicon chips independently run the variables through an HMAC-SHA1 mathematical permutation pipeline simultaneously.
4. The immense resulting matrix is dynamically truncated and bitwise-shifted into a globally synchronized 6-digit numeral.

Even possessing the primary administrative password, remote attackers are crippled without concurrent possession of the physical phone displaying the temporal integer. 

% -------------------------
% NODE MCU LEVEL 10
% -------------------------
\section{Systematic NodeMCU Implementation (Demo 10)}
A standard microcontroller lacks the raw 64-bit integer space necessary to easily crunch massive encryption algorithms. Fortunately, the Espressif ESP8266 core encompasses the \texttt{BearSSLHelpers} cryptography suite embedded deep beneath the network stack.

\begin{lstlisting}[language=C++, caption=Demo\_10\_Secure\_IoT\_Admin\_TOTP\_With\_Description.ino]
/*
  ===============================================================
  Secure IoT System - Admin + Password + TOTP (QR + Manual)
  ===============================================================

  FEATURES:
  ✓ Admin login (username: admin, password: 1234)
  ✓ Google Authenticator TOTP (6-digit, 30-second)
  ✓ QR code generation using api.qrserver.com
  ✓ Manual secret setup option
  ✓ Auto-refresh live sensor display (every 5 sec)
  ✓ Secure /set route (requires password + fresh OTP)
  ✓ 3 wrong update attempts -> "Unauthorized access"
  ✓ Cyan temperature, Pink humidity display

  REQUIRED LIBRARIES:
  - Base32-Decode
  - NTPClient
  (BearSSL is included in ESP8266 core)

  HARDWARE:
  - NodeMCU (ESP8266)
  - DHT22 connected to D5
*/

#include <ESP8266WiFi.h>
#include <ESP8266WebServer.h>
#include <DHT.h>
#include <WiFiUdp.h>
#include <NTPClient.h>
#include <Base32-Decode.h>
#include <BearSSLHelpers.h>

const char* ssid = "drarka.in";
const char* password = "arkaprabha@1985";

const String ADMIN_USER = "admin";
const String ADMIN_PASS = "1234";

const String SECRET_KEY = "JBSWY3DPEHPK3PXPJBSWY3DPEHPK3PXP";

#define DHTPIN D5
#define DHTTYPE DHT22
DHT dht(DHTPIN, DHTTYPE);

ESP8266WebServer server(80);
WiFiUDP ntpUDP;
NTPClient timeClient(ntpUDP, "pool.ntp.org", 0, 60000);

uint8_t decodedKey[32];
int decodedKeyLen = 0;

bool loggedIn = false;
int failedAttempts = 0;
bool blocked = false;

float overrideTemp = NAN;
float overrideHum  = NAN;

// Step 1: The Core HMAC-SHA1 Algorithm translated to C++ Pointers
String getTOTP(unsigned long epochSeconds) {

  uint64_t counter = epochSeconds / 30;

  uint8_t msg[8];
  for (int i = 7; i >= 0; i--) {
    msg[i] = counter & 0xFF;
    counter >>= 8;
  }

  uint8_t hash[20];
  br_hmac_key_context kc;
  br_hmac_context ctx;

  br_hmac_key_init(&kc, &br_sha1_vtable, decodedKey, decodedKeyLen);
  br_hmac_init(&ctx, &kc, 0);
  br_hmac_update(&ctx, msg, sizeof(msg));
  br_hmac_out(&ctx, hash);

  int offset = hash[19] & 0x0F;

  uint32_t binary =
    ((hash[offset] & 0x7F) << 24) |
    ((hash[offset+1] & 0xFF) << 16) |
    ((hash[offset+2] & 0xFF) << 8) |
    (hash[offset+3] & 0xFF);

  uint32_t otp = binary % 1000000;

  char buf[7];
  snprintf(buf, sizeof(buf), "%06u", otp);

  return String(buf);
}

void handleRoot() {

  if (!loggedIn) {

    if (!server.hasArg("user")) {
      server.send(200, "text/html",
        "<h3>Admin Login</h3>"
        "<form method='GET'>"
        "Username:<br><input name='user'><br>"
        "Password:<br><input name='pass' type='password'><br><br>"
        "<input type='submit' value='Login'>"
        "</form>");
      return;
    }

    if (server.arg("user") != ADMIN_USER ||
        server.arg("pass") != ADMIN_PASS) {
      server.send(401, "text/plain", "Invalid credentials");
      return;
    }

    if (!server.hasArg("otp")) {

      String otpURL =
        "otpauth://totp/IoT-Secure?secret=" +
        SECRET_KEY +
        "&issuer=IoT-Secure";

      String qrURL =
        "https://api.qrserver.com/v1/create-qr-code/?size=300x300&data=" +
        otpURL;

      String html =
        "<h3>Google Authenticator Setup</h3>"
        "<img src='" + qrURL + "' width='250'><br><br>"
        "<p><b>Manual Setup Key:</b> " + SECRET_KEY + "</p>"
        "<hr>"
        "<form method='GET'>"
        "<input type='hidden' name='user' value='admin'>"
        "<input type='hidden' name='pass' value='1234'>"
        "Enter 6-digit OTP:<br><input name='otp'><br><br>"
        "<input type='submit' value='Verify'>"
        "</form>";

      server.send(200, "text/html", html);
      return;
    }

    timeClient.update();

    if (server.arg("otp") == getTOTP(timeClient.getEpochTime())) {
      loggedIn = true;
    } else {
      server.send(401, "text/plain", "Invalid OTP");
      return;
    }
  }

  float t = isnan(overrideTemp) ? dht.readTemperature() : overrideTemp;
  float h = isnan(overrideHum)  ? dht.readHumidity()    : overrideHum;

  String html =
    "<meta http-equiv='refresh' content='5'>"
    "<body style='background:black;color:white;text-align:center;'>"
    "<h2>Secure Live Sensor Data</h2>"
    "<h3>Temperature</h3>"
    "<h1 style='color:cyan;'>" + String(t,1) + " C</h1>"
    "<h3>Humidity</h3>"
    "<h1 style='color:pink;'>" + String(h,1) + " %</h1>"
    "</body>";

  server.send(200, "text/html", html);
}

void handleSet() {

  if (blocked) {
    server.send(403, "text/plain", "Unauthorized access");
    return;
  }

  if (!server.hasArg("pass") || !server.hasArg("otp")) {
    server.send(400, "text/plain", "Password and OTP required");
    return;
  }

  timeClient.update();

  if (server.arg("pass") != ADMIN_PASS ||
      server.arg("otp") != getTOTP(timeClient.getEpochTime())) {

    failedAttempts++;

    if (failedAttempts >= 3) {
      blocked = true;
      server.send(403, "text/plain", "Unauthorized access");
      return;
    }

    server.send(401, "text/plain", "Invalid credentials");
    return;
  }

  failedAttempts = 0;

  if (server.hasArg("t")) overrideTemp = server.arg("t").toFloat();
  if (server.hasArg("h")) overrideHum  = server.arg("h").toFloat();

  server.send(200, "text/plain", "Values updated successfully");
}

void setup() {

  Serial.begin(115200);
  dht.begin();

  WiFi.begin(ssid, password);
  while (WiFi.status() != WL_CONNECTED) delay(300);

  Serial.println("Connected!");
  Serial.println(WiFi.localIP());

  timeClient.begin();
  while (!timeClient.update()) timeClient.forceUpdate();

  decodedKeyLen = base32decode(
    SECRET_KEY.c_str(),
    decodedKey,
    sizeof(decodedKey)
  );

  server.on("/", handleRoot);
  server.on("/set", handleSet);
  server.begin();
}

void loop() {
  server.handleClient();
  timeClient.update();
}
\end{lstlisting}

% -------------------------
% PICO W AUTHENTICATOR
% -------------------------
\section{Pico W Level 10: Official Authenticator Upgrade}
The MicroPython developer finally merges universal Y2K Unix Epoch mathematical fixes and the official `hmac-sha1` integer truncation sequence, completely matching Google Authenticator specifications.

\begin{lstlisting}[language=Python, caption=07\_strong\_login\_Authenticator.py]
import network, socket, time, machine, dht, ubinascii, secrets, hashlib, rp2, ntptime

# --- [1. WIFI & TIME SYNC] ---
rp2.country('US')
sensor = dht.DHT22(machine.Pin(15))
wlan = network.WLAN(network.STA_IF)
wlan.active(True)
wlan.config(pm = 0xa11140) 
wlan.connect(secrets.SSID, secrets.PASSWORD)

print("Connecting to WiFi", end="")
while wlan.status() < 3:
    print(".", end="")
    time.sleep(1)

my_ip = wlan.ifconfig()[0]

print("\nSyncing exact time with Internet (NTP)", end="")
ntp_success = False
for i in range(10): # Try up to 10 times to ensure we get the real time
    try:
        ntptime.settime()
        ntp_success = True
        break
    except:
        time.sleep(1)
        print(".", end="")

if ntp_success:
    print(" [SUCCESS] Clock synced!")
else:
    print("\n[!!!] CLOCK SYNC FAILED [!!!] Codes will not match.")

# --- [2. CRITICAL FIX: EPOCH DETECTION] ---
if time.localtime(0)[0] == 2000:
    UNIX_OFFSET = 946684800
    print("[SYSTEM] Y2K Epoch detected. Applying 30-year correction.")
else:
    UNIX_OFFSET = 0

# --- [3. BULLETPROOF GOOGLE AUTHENTICATOR SETUP] ---
print("\n" + "="*50)
print("  ADD THIS NEW KEY TO GOOGLE AUTHENTICATOR")
print("="*50)
print("Option 1: Scan this QR Code from your browser:")
print("https://api.qrserver.com/v1/create-qr-code/?data=otpauth://totp/Pico:Pico%20Monitor?secret=GEZDGNBVGY3TQOJQGEZDGNBVGY3TQOJQ&issuer=Pico&size=300x300")
print("Option 2: Enter Key Manually:")
print("Account: Pico Monitor")
print("Key:     GEZDGNBVGY3TQOJQGEZDGNBVGY3TQOJQ")
print("-" * 50)
print(f"[LIVE DASHBOARD] http://{my_ip}")
print("="*50 + "\n")

auth = ubinascii.b2a_base64(f"{secrets.WEB_USER}:{secrets.WEB_PASS_STRONG}".encode()).strip().decode()
expected_basic = f"Basic {auth}"

TOTP_INTERVAL = 30 

SECRET_BYTES = b"12345678901234567890"

authorized_ips = {}
fail_counts = {}
MAX_FAILS = 3
last_totp, last_sec = "", -1

def hmac_sha1(key, msg):
    """Official HMAC-SHA1 algorithm"""
    if len(key) > 64: key = hashlib.sha1(key).digest()
    key = key + b'\x00' * (64 - len(key))
    ipad = bytes(x ^ 0x36 for x in key)
    opad = bytes(x ^ 0x5C for x in key)
    inner = hashlib.sha1(ipad + msg).digest()
    return hashlib.sha1(opad + inner).digest()

def get_google_auth_code():
    """Generates the 6-digit Google Auth code using real UTC Time"""
    unix_time = time.time() + UNIX_OFFSET
    time_step = int(unix_time // TOTP_INTERVAL)
    
    msg = bytearray(8)
    for i in range(7, -1, -1):
        msg[i] = time_step & 0xff
        time_step >>= 8
        
    hash_result = hmac_sha1(SECRET_BYTES, msg)
    offset = hash_result[19] & 0x0F
    code_int = ((hash_result[offset] & 0x7F) << 24 |
                (hash_result[offset+1] & 0xFF) << 16 |
                (hash_result[offset+2] & 0xFF) << 8 |
                (hash_result[offset+3] & 0xFF))
    
    code = str(code_int % 1000000)
    return "0" * (6 - len(code)) + code

# --- [4. SOCKET ENGINE] ---
s = socket.socket()
s.setsockopt(socket.SOL_SOCKET, socket.SO_REUSEADDR, 1)
s.bind(('0.0.0.0', 80))
s.listen(1)
s.settimeout(0.1)

t, h, spoof = "0.0", "0.0", False

# --- [5. MAIN LOOP] ---
try:
    while True:
        curr_t = int(time.time())
        code = get_google_auth_code()
        
        if curr_t != last_sec:
            last_sec = curr_t
            if code != last_totp: last_totp = code
            print(f"\r[Pico Code: {last_totp}] | Matches phone? Rotates in: {TOTP_INTERVAL-(curr_t%TOTP_INTERVAL):02d}s ", end="")

        if not spoof:
            try: 
                sensor.measure()
                t, h = "{:.1f}".format(sensor.temperature()), "{:.1f}".format(sensor.humidity())
            except: pass

        conn = None
        try:
            conn, addr = s.accept()
            ip = addr[0]
            conn.settimeout(0.5)
            
            req = conn.recv(1024).decode()
            if not req:
                conn.close()
                continue

            if ip in fail_counts and fail_counts[ip] >= MAX_FAILS:
                conn.sendall(b'HTTP/1.1 403 Forbidden\r\nConnection: close\r\n\r\n<h1>IP BLOCKED: Too Many Failed 2FA Attempts</h1>')
                conn.close()
                continue

            if expected_basic in req:
                
                if 'code=' in req:
                    if f"code={last_totp}" in req:
                        authorized_ips[ip] = time.time() + 600 # 10 minute session
                        print(f"\n[AUTH SUCCESS] {ip} unlocked dashboard.")
                        conn.sendall(b'HTTP/1.1 302 Found\r\nLocation: /\r\nConnection: close\r\n\r\n')
                        conn.close()
                        continue
                    else:
                        fail_counts[ip] = fail_counts.get(ip, 0) + 1
                        print(f"\n[AUTH FAILED] {ip} Wrong Code. ({fail_counts[ip]}/{MAX_FAILS})")
                        conn.sendall(b'HTTP/1.1 403 Forbidden\r\nConnection: close\r\n\r\n<h2>Invalid Google Auth Code</h2>')
                        conn.close()
                        continue

                if ip in authorized_ips and time.time() < authorized_ips[ip]:
                    html = f"""<!DOCTYPE html><html><head><meta http-equiv="refresh" content="2;url=/">
                    <style>
                        body {{ background:#000; color:white; font-family:sans-serif; display:flex; flex-direction:column; align-items:center; justify-content:center; height:100vh; margin:0; }}
                        .l {{ color:#333; font-size:1.5rem; letter-spacing:5px; margin-top:20px; font-weight:bold; }}
                        .t {{ color:cyan; font-size:10rem; font-weight:bold; margin:0; text-shadow:0 0 30px cyan; }}
                        .h {{ color:pink; font-size:10rem; font-weight:bold; margin:0; text-shadow:0 0 30px pink; }}
                    </style></head><body>
                        <div class="l">TEMPERATURE</div><div class="t">{t}&deg;</div>
                        <div class="l">HUMIDITY</div><div class="h">{h}%</div>
                    </body></html>"""
                else:
                    html = f"""<!DOCTYPE html><html><head><meta name="viewport" content="width=device-width, initial-scale=1">
                    <style>
                        body {{ background:#0a0a0a; color:cyan; font-family:sans-serif; text-align:center; padding-top:20vh; margin:0; }}
                        input {{ background:black; border:2px solid cyan; color:cyan; padding:15px; font-size:1.5rem; border-radius:10px; width:200px; text-align:center; outline:none; letter-spacing:3px; }}
                        button {{ background:cyan; border:none; padding:15px 30px; font-size:1rem; font-weight:bold; border-radius:10px; margin-top:20px; color:black; cursor:pointer; }}
                    </style></head><body>
                        <h2>GOOGLE AUTHENTICATOR</h2>
                        <form action="/" method="get">
                            <input type="text" name="code" placeholder="000 000" inputmode="numeric" autofocus autocomplete="off"><br>
                            <button type="submit">UNLOCK DASHBOARD</button>
                        </form></body></html>"""
                
                response = f'HTTP/1.1 200 OK\r\nContent-Type: text/html\r\nConnection: close\r\n\r\n{html}'
                conn.sendall(response.encode('utf-8'))
            else:
                conn.sendall(b'HTTP/1.1 401 Unauthorized\r\nWWW-Authenticate: Basic realm="Pico"\r\nConnection: close\r\n\r\n')
            
            conn.close()
        except OSError:
            pass
        
        time.sleep(0.02)
except KeyboardInterrupt:
    print("\n[SHUTDOWN] Stopping Server...")
finally:
    s.close()
\end{lstlisting}

\begin{takeawaybox}
By enforcing a strict HMAC-SHA1 algorithm mathematically matching the global RFC standard, the devices are officially fully secure. A hacker intercepting the `Basic` hash and circumventing the 3-strike network lockout will still fail against the Time-based temporal shifting properties of independent local execution arrays. You have secured the machine!
\end{takeawaybox}
